% GENERATED FILE. Edit canonical lessons, then run npm run generate:books.
% canonical-source-hash: fnv1a64:d4ead0fc11c63ab9
% canonical-lessons: TA-C260-atikari2, TA-C260-paravu, TA-C260-mulku, TA-C260-mita, TA-C260-pili

\chapter{To increase, To spread, To sink, To float, To squeeze}
\label{ch:ta-to-increase-to-spread-to-sink-to-float-to-squeeze}
\hlchaptermodality{260}

\begin{chapteropening}
\textbf{By the end of this chapter:} \emph{I can say to increase, to spread, to sink, to float and to squeeze.}

Complete the last lesson of chapter 260: I can say to increase, to spread, to sink, to float and to squeeze.
\end{chapteropening}

\section[\texorpdfstring{\ta{அதிகரி} (atikari)}{அதிகரி (atikari)}]{\ta{அதிகரி} (atikari) — to increase, to go up}

\label{lesson:TA-C260-atikari2}



\begin{quote}
Before the new one: say the Tamil for to repair, then the Tamil for to tidy up.
\end{quote}

\subsection*{You'll want to know: \ta{அதிகரி}}
\textbf{\ta{அதிகரி}} — \emph{atikari} — \textquotedblleft{}to increase, to go up\textquotedblright{}.

\begin{grammarlens}[title={What you've built}]
One more word for this chapter.
\end{grammarlens}

\subsection*{Your turn}
\begin{itemize}
\raggedright
  \item \emph{Say it:} \emph{atikari}
  \item \emph{Say it:} \emph{atikari}, once more
  \item \emph{Say it:} say \emph{atikari}
  \item \emph{Recall:} say the Tamil for to repair, then the Tamil for to tidy up, then say \emph{atikari} again
\end{itemize}

\begin{tcolorbox}[breakable,colback=teal!4,colframe=teal!35!black,title={Before you move on}]
What does \ta{அதிகரி} mean? (\textquotedblleft{}To increase, to go up\textquotedblright{}.) Say it once more.
\end{tcolorbox}
\section[\texorpdfstring{\ta{பரவு} (paravu)}{பரவு (paravu)}]{\ta{பரவு} (paravu) — to spread (of news, disease, fire)}

\label{lesson:TA-C260-paravu}



\begin{quote}
Before the new one: say the Tamil for to tidy up, then the Tamil for to increase.
\end{quote}

\subsection*{You'll want to know: \ta{பரவு}}
\textbf{\ta{பரவு}} — \emph{paravu} — \textquotedblleft{}to spread (of news, disease, fire)\textquotedblright{}.

\begin{grammarlens}[title={What you've built}]
One more word for this chapter.
\end{grammarlens}

\subsection*{Your turn}
\begin{itemize}
\raggedright
  \item \emph{Say it:} \emph{paravu}
  \item \emph{Say it:} \emph{paravu}, once more
  \item \emph{Say it:} say \emph{paravu}
  \item \emph{Recall:} say the Tamil for to tidy up, then the Tamil for to increase, then say \emph{paravu} again
\end{itemize}

\begin{tcolorbox}[breakable,colback=teal!4,colframe=teal!35!black,title={Before you move on}]
What does \ta{பரவு} mean? (\textquotedblleft{}To spread (of news, disease, fire)\textquotedblright{}.) Say it once more.
\end{tcolorbox}
\section[\texorpdfstring{\ta{மூழ்கு} (mūḻku)}{மூழ்கு (mūḻku)}]{\ta{மூழ்கு} (mūḻku) — to sink, to drown}

\label{lesson:TA-C260-mulku}



\begin{quote}
Before the new one: say the Tamil for to increase, then the Tamil for to spread.
\end{quote}

\subsection*{You'll want to know: \ta{மூழ்கு}}
\textbf{\ta{மூழ்கு}} — \emph{mūḻku} — \textquotedblleft{}to sink, to drown\textquotedblright{}.

\begin{grammarlens}[title={What you've built}]
One more word for this chapter.
\end{grammarlens}

\subsection*{Your turn}
\begin{itemize}
\raggedright
  \item \emph{Say it:} \emph{mūḻku}
  \item \emph{Say it:} \emph{mūḻku}, once more
  \item \emph{Say it:} say \emph{mūḻku}
  \item \emph{Recall:} say the Tamil for to increase, then the Tamil for to spread, then say \emph{mūḻku} again
\end{itemize}

\begin{tcolorbox}[breakable,colback=teal!4,colframe=teal!35!black,title={Before you move on}]
What does \ta{மூழ்கு} mean? (\textquotedblleft{}To sink, to drown\textquotedblright{}.) Say it once more.
\end{tcolorbox}
\section[\texorpdfstring{\ta{மித} (mita)}{மித (mita)}]{\ta{மித} (mita) — to float}

\label{lesson:TA-C260-mita}



\begin{quote}
Before the new one: say the Tamil for to spread, then the Tamil for to sink.
\end{quote}

\subsection*{You'll want to know: \ta{மித}}
\textbf{\ta{மித}} — \emph{mita} — \textquotedblleft{}to float\textquotedblright{}.

\begin{grammarlens}[title={What you've built}]
One more word for this chapter.
\end{grammarlens}

\subsection*{Your turn}
\begin{itemize}
\raggedright
  \item \emph{Say it:} \emph{mita}
  \item \emph{Say it:} \emph{mita}, once more
  \item \emph{Say it:} say \emph{mita}
  \item \emph{Recall:} say the Tamil for to spread, then the Tamil for to sink, then say \emph{mita} again
\end{itemize}

\begin{tcolorbox}[breakable,colback=teal!4,colframe=teal!35!black,title={Before you move on}]
What does \ta{மித} mean? (\textquotedblleft{}To float\textquotedblright{}.) Say it once more.
\end{tcolorbox}
\section[\texorpdfstring{\ta{பிழி} (piḻi)}{பிழி (piḻi)}]{\ta{பிழி} (piḻi) — to squeeze, to wring out}

\label{lesson:TA-C260-pili}



\begin{quote}
Before the new one: say the Tamil for to sink, then the Tamil for to float.
\end{quote}

\subsection*{You'll want to know: \ta{பிழி}}
\textbf{\ta{பிழி}} — \emph{piḻi} — \textquotedblleft{}to squeeze, to wring out\textquotedblright{}.

\begin{grammarlens}[title={What you've built}]
One more word for this chapter.
\end{grammarlens}

\subsection*{Your turn}
\begin{itemize}
\raggedright
  \item \emph{Say it:} \emph{piḻi}
  \item \emph{Say it:} \emph{piḻi}, once more
  \item \emph{Say it:} say \emph{piḻi}
  \item \emph{Recall:} say the Tamil for to sink, then the Tamil for to float, then say \emph{piḻi} again
\end{itemize}

\begin{tcolorbox}[breakable,colback=teal!4,colframe=teal!35!black,title={Before you move on}]
What does \ta{பிழி} mean? (\textquotedblleft{}To squeeze, to wring out\textquotedblright{}.) Say it once more.
\end{tcolorbox}
